\section{点态与一致收敛只差一个量词顺序}
\label{sec:12-Real-Analysis-Support/02-Convergence/02}

你在逐步改进一个函数逼近器，每训练一轮就得到一条新曲线。挑任何一个输入点盯着看，它的误差确实一路降到零；可无论跑到第几轮，总能在输入空间里找到一小块区域，那里的误差还是很大——只不过这块区域每一轮都换个位置。逐点的好消息和整体的坏消息同时为真，两者并不矛盾。

上一节的 \(\varepsilon\)--\(N\) 定义里，\(N\) 只依赖 \(\varepsilon\)。函数列多了一个自变量，于是 \(N\) 多了一个可以依赖的对象：点 \(x\)。允不允许这份依赖，正是点态收敛与一致收敛的全部分别。

\begin{center}
\begin{tikzpicture}[x=5.2cm,y=3.1cm,>=stealth]
  \draw[thick,->] (0,0) -- (1.2,0) node[right] {\footnotesize \(x\)};
  \draw[thick,->] (0,0) -- (0,1.25);
  \draw (1,0.025) -- (1,-0.025) node[below] {\footnotesize \(1\)};
  \draw (0.02,1) -- (-0.02,1) node[left] {\footnotesize \(1\)};
  \draw[dashed] (0,0.25) -- (1.06,0.25);
  \node[left] at (0,0.25) {\footnotesize \(\varepsilon\)};
  \foreach \n in {1,2,4,10,40}
    \draw[gray] plot[domain=0:1,samples=140] (\x,{pow(\x,\n)});
  \draw[very thick] (0,0) -- (1,0);
  \draw[thick,fill=white] (1,0) circle (1.7pt);
  \fill (1,1) circle (1.7pt);
  \draw[->] (0.42,0.62) -- (0.78,0.22);
  \node[above left] at (0.44,0.6) {\footnotesize \(n\) 增大};
  \draw[->] (0.70,0.92) -- (0.955,0.58);
  \node[align=center,above] at (0.66,0.9) {\footnotesize 每条曲线都在\\\footnotesize 端点前陡起};
\end{tikzpicture}
\end{center}

\emph{\(f_n(x)=x^n\) 越来越贴着底边，极限是一条折线加一个孤立的端点值；超出 \(\varepsilon\) 带的坏区域向 \(1\) 收缩，却从不消失。}

\subsection{两个定义只差 \(\forall x\) 与 \(\exists N\) 的先后}

设 \(f_n:E\to\R\)。点态收敛到 \(f\) 写成

\[
\forall x\in E,\ \forall\varepsilon>0,\ \exists N=N(x,\varepsilon),\ \forall n\ge N,\ \abs{f_n(x)-f(x)}<\varepsilon,
\]

一致收敛写成

\[
\forall\varepsilon>0,\ \exists N=N(\varepsilon),\ \forall x\in E,\ \forall n\ge N,\ \abs{f_n(x)-f(x)}<\varepsilon.
\]

两串符号的字母完全相同，只是 \(\exists N\) 从 \(\forall x\) 的后面挪到了前面。挪之前，每个点自带一张进度表，你可以先看点再定时刻；挪之后，你必须先交出一个时刻，它得对全体点同时兑现。一致收敛蕴含点态收敛，反向不成立——差的那部分正是本节要展开的东西。

第三种写法最便于计算：一致收敛等价于

\[
\norm{f_n-f}_\infty = \sup_{x\in E}\abs{f_n(x)-f(x)} \longrightarrow 0.
\]

取上确界这个动作就是在替所有点统一发言。有了它，函数列的收敛被压成一个普通数列的收敛，上一章关于数列的全部工具立刻可用。

\subsection{\(x^n\) 把差别演成一场追不上的追逐}

在 \(E=[0,1]\) 上取 \(f_n(x)=x^n\)。固定 \(x<1\) 时 \(x^n\to0\)，而 \(f_n(1)=1\) 恒成立，所以点态极限是

\[
f(x)=
\begin{cases}
0, & 0\le x<1,\\[2pt]
1, & x=1.
\end{cases}
\]

每个 \(f_n\) 都是多项式，处处连续；极限却在 \(x=1\) 处断开。连续性没有被继承下来。

不一致的原因图上看得很清楚。固定 \(n\)，让 \(x\) 从左边趋近 1，由连续性 \(x^n\) 可以任意接近 1，于是

\[
\sup_{x\in[0,1]}\abs{f_n(x)-f(x)} = 1
\]

对每个 \(n\) 都成立，sup 范数根本没有下降。换成动态的说法：对每个 \(n\)，右端附近有一块坏区域，那里 \(x^n\) 还没掉进 \(\varepsilon\) 带；\(n\) 增大时坏区域向 1 收缩，但永远保有正的宽度。任何固定的 \(x<1\) 终会脱离坏区域——不是因为那块区域被修好了，而是因为它挪到了 \(x\) 的右边去。开篇那个每轮换位置的高误差区域，是同一件事的工程版本。

定义域一改，结论就翻转。在 \([0,a]\) 上（\(a<1\)），

\[
\sup_{x\in[0,a]} x^n = a^n \to 0,
\]

一致收敛成立，因为坏区域被推到了区间外面。一致性从来不是函数列单方面的属性，它是函数列与定义域一起决定的。这也解释了为什么模型在核心数据分布上表现稳定、一到边界区域就崩——问题往往不在逼近能力，而在于评估用的那个域是否包含了坏区域盘踞的地方。

\subsection{保住连续性靠的是``先定 \(N\)，再定 \(\delta\)''}

\begin{theorem}
若每个 \(f_n\) 在 \(E\) 上连续，且 \(f_n\to f\) 一致，则 \(f\) 在 \(E\) 上连续。
\end{theorem}

固定 \(x_0\)，把误差拆成三段：

\[
\begin{aligned}
\abs{f(x)-f(x_0)} \le\ & \abs{f(x)-f_N(x)}\\
& + \abs{f_N(x)-f_N(x_0)}\\
& + \abs{f_N(x_0)-f(x_0)}.
\end{aligned}
\]

给定 \(\varepsilon>0\)，先用一致性挑出一个 \(N\)，使第一、三项对所有 \(x\) 都小于 \(\varepsilon/3\)；\(N\) 就此钉死。再用 \(f_N\) 这一个函数的连续性挑出 \(\delta>0\)，使 \(\abs{x-x_0}<\delta\) 时中间项小于 \(\varepsilon/3\)。三段合起来小于 \(\varepsilon\)。

条件对账在于顺序。第一步需要一个不依赖 \(x\) 的 \(N\)，否则第二步就没有确定的 \(f_N\) 可以调用其连续性；而点态收敛只能给出 \(N(x,\varepsilon)\)，你想固定它，它就随着 \(x\) 变。整个证明卡在第一步，卡的位置与 \(x^n\) 的坏区域是同一处。一致收敛买下的，正是``先全局选 \(N\)、再局部选 \(\delta\)''这个操作次序。

同样的次序也让极限穿过积分号。在有限区间 \([a,b]\) 上，

\[
\abs{\int_a^b f_n - \int_a^b f} \le (b-a)\,\norm{f_n-f}_\infty \to 0,
\]

一行不等式就够了：一致的误差界乘以区间长度，仍然趋于零。点态收敛给不出这个乘法所需的统一因子，\hyperref[sec:12-Real-Analysis-Support/04-Limits-and-Integration/01]{``点态极限为何不能直接穿过积分号''}给出了积分与极限交换失败的具体样本。

\subsection{sup 范数把函数列变回数列}

有了 \(\norm{\cdot}_\infty\)，Cauchy 判据可以照搬。称 \((f_n)\) 一致 Cauchy，若对任意 \(\varepsilon>0\) 存在 \(N\) 使 \(m,n\ge N\) 时 \(\sup_{x\in E}\abs{f_m(x)-f_n(x)}<\varepsilon\)。从这个条件造出极限函数分两步：固定 \(x\)，\((f_n(x))\) 是 \(\R\) 中的 Cauchy 列，实数完备性给出一个数，记作 \(f(x)\)，这样先得到极限函数的逐点定义；再在不等式中令 \(m\to\infty\)，得 \(\abs{f_n(x)-f(x)}\le\varepsilon\) 对所有 \(x\) 与 \(n\ge N\) 成立，右端与 \(x\) 无关，取上确界即得一致收敛。

这条路径的好处与上一章的 Cauchy 判据一致：不必先猜出 \(f\) 长什么样。紧集 \(K\) 上的连续函数空间 \(C(K)\) 配 sup 范数是完备的，连续函数的一致 Cauchy 列收敛到连续函数，极限不会跑出空间——这正是 Banach 空间这个名字所承诺的内容，也是 \hyperref[sec:12-Real-Analysis-Support/03-Topology-Basics/03]{``完备性让逼近留在空间内''}会重新拾起的话题。

函数级数用起来尤其顺手。若 \(\sum_n\sup_{x\in E}\abs{g_n(x)}<\infty\)，则部分和满足

\[
\sup_{x\in E}\abs{S_m(x)-S_n(x)} \le \sum_{k=n+1}^{m}\ \sup_{x\in E}\abs{g_k(x)},
\]

右端是一个收敛正项级数的尾部，趋于零。于是部分和一致 Cauchy，级数一致收敛。这就是 Weierstrass 判别法的全部内容——正项级数的比较判别法换了身衣服。

\subsection{一致收敛管函数值，不管别的}

它的势力范围有明确边界。取 \(f_n(x)=\frac{\sin(nx)}{n}\)，由 \(\abs{f_n}\le1/n\) 知它在 \(\R\) 上一致收敛到零；可导数 \(f_n'(x)=\cos(nx)\) 在每一点都在 \(-1\) 与 \(1\) 之间振荡，连点态极限都没有。函数值被压平了，变化率没有。想交换极限与求导，得对导数列另立一致收敛的条件。

定义域无限时，积分的交换也会失守。\(\R\) 上的常函数 \(f_n\equiv 1/n\) 一致收敛到零，但每个 \(f_n\) 的积分都是无穷，极限函数的积分是零。前面那行不等式的因子 \(b-a\) 在这里变成了无穷大，小误差乘以无限大的域，乘积失控。

一致收敛因此是一档具体的强度，不是万能通行证：它足以保住连续性，足以在有限区间上换积分，不足以换导数，也不足以对付无限测度的定义域。要交换的运算越凶，对``接近''的要求越高。接下来要问的正是这个问题——\hyperref[ch:100]{``什么时候可以交换极限与运算''}把各类条件按运算逐一配平，控制收敛定理会用一个可积包络替下这里的上确界。
